\begin{textsample}{NEG Dim 1 – vem\_ed\_25 – Score -2.00 – vem\_ed\_25/t131.txt}  \label{ex:f1_neg_116}
Leitores Aos Osvaldo Succi Jr.
Coordenador PCIs Esta edição reflete a \textbf{movimentada} \textbf{agenda} de eventos acadêmicos do segundo semestre de 2024.
Em 29 de \textbf{agosto}, ocorreu a 3ª Jornada RedAES, Rede de Apoio ao Ensino Superior formada por nove instituições paulistas: Centro Paula Souza, IFSP, UFABC, UFSCar, Unesp, Unicamp, Unifesp, Univesp e USP.
On-line e gratuita, a Jornada apresentou três mesas temáticas, que \textbf{somaram} 3.954 participantes.
Tive o prazer de \textbf{moderar} a mesa 2, dedicada ao debate das tendências futuras dos Intercâmbios Virtuais.
O tema do 4º Congresso da Red LatAM COIL foi “diversidade, equidade e inclusão”.
Ana Carolina Freschi, da nossa equipe de PCIs / Cesu, integra o conselho executivo da rede e participou da abertura do evento, em 4 de setembro.
No dia 5, apresentei workshop com os colegas José Luis Jiménez (Universidad Católica Andrés Bello, Venezuela) e Mirjam Hauck (The Open University, Reino Unido).
No 5º Congresso de Línguas para Fins Específicos (LinFE), Ana Carolina Freschi, Priscilla Ferro e Regiane Camargo Moreira (equipe dos PCIs / Cesu) compartilharam a experiência dos Projetos Colaborativos Internacionais no Centro Paula Souza.
As professoras Patrícia Januária Barbosa e Tálita Guarino também apresentaram trabalho, sobre PCI realizado na Fatec Guaratinguetá.
Na seção “Artigo de Opinião”, Edilene Gasparini Fernandes comenta sobre o uso das WebQuests nos PCIs que orienta na Fatec São José do Rio Preto.
Para \textbf{enviar} seu artigo, acesse https://publicacoescesu.cps.sp.gov.br/VEm/about/submissions.
Dúvidas?
Escreva para cesu.pci@fatec.sp.gov.br Boa leitura!

% matched lemmas: agenda, agosto, enviar, moderar, movimentado, somar
\end{textsample}
